%=====================================================================
% Datos · Casos que la base rechaza
%=====================================================================
\subsection*{Casos que la base rechaza}
\addcontentsline{toc}{subsection}{Casos que la base rechaza}

Unos datos de prueba solo sirven si la base también rechaza los que no deberían existir. \at{bd/pruebas\_rechazo.sql} intenta, sobre los registros cargados, \dnNumPruebasRechazo{} operaciones que el documento prohíbe, cada una en su propia transacción, que se revierte siempre: el script no cambia ningún dato. De cada intento anota el número de error de SQL Server y la restricción que lo impidió, que sale del propio mensaje de error.

% Archivo generado por bd/exportar_datos.py. No editar a mano.
\begin{center}\footnotesize\setlength{\tabcolsep}{2pt}
\begin{longtable}{|L{29.5pt}|L{166.5pt}|L{83.0pt}|L{26.1pt}|L{143.9pt}|}
\hline
\textbf{\hspace{0pt}Núm.} & \textbf{\hspace{0pt}Operación que se intenta} & \textbf{\hspace{0pt}Regla} & \textbf{\hspace{0pt}Error} & \textbf{\hspace{0pt}Rechazada por} \\ \hline
\endfirsthead
\hline
\textbf{\hspace{0pt}Núm.} & \textbf{\hspace{0pt}Operación que se intenta} & \textbf{\hspace{0pt}Regla} & \textbf{\hspace{0pt}Error} & \textbf{\hspace{0pt}Rechazada por} \\ \hline
\endhead
1 & Estado de cuenta fuera del dominio & CRUD, dominios & 547 & CK\_\allowbreak{}Usuario\_\allowbreak{}estado\_\allowbreak{}cuenta \\ \hline
2 & Nickname igual salvo mayúsculas y acentos & CU-02 RN-01 & 2627 & UQ\_\allowbreak{}Usuario\_\allowbreak{}nickname \\ \hline
3 & Idioma no admitido & CU-12 FA-07, D-21 & 547 & CK\_\allowbreak{}Usuario\_\allowbreak{}idioma \\ \hline
4 & Saldo de monedas negativo & CU-40 RN-02 & 547 & CK\_\allowbreak{}Usuario\_\allowbreak{}rangos \\ \hline
5 & Amistad con una cuenta que no existe & Llave foránea & 547 & FK\_\allowbreak{}Amistad\_\allowbreak{}UsuarioB \\ \hline
6 & Equipar un objeto que no se posee & CU-14 RN-02 & 547 & FK\_\allowbreak{}Equipamiento\_\allowbreak{}UsuarioObjeto \\ \hline
7 & Equipar un objeto en otra ranura & D-03, CU-14 RN-01 & 547 & FK\_\allowbreak{}Equipamiento\_\allowbreak{}ObjetoCosmetico \\ \hline
8 & Segunda partida sin terminar & D-11 & 2601 & UX\_\allowbreak{}Participacion\_\allowbreak{}sin\_\allowbreak{}terminar \\ \hline
9 & Segundo código de recuperación vigente & CU-03 RN-05 & 2601 & UX\_\allowbreak{}TokenRecuperacion\_\allowbreak{}vigente \\ \hline
10 & Reportarse a sí mismo & CU-42 FA-06 & 547 & CK\_\allowbreak{}Reporte\_\allowbreak{}distintos \\ \hline
11 & Sanción temporal sin fecha de fin & CU-44 RN-04 & 547 & CK\_\allowbreak{}Sancion\_\allowbreak{}tipo \\ \hline
12 & Compra que no dice qué objeto & CU-38 RN-08 & 547 & CK\_\allowbreak{}MovimientoMoneda\_\allowbreak{}origen \\ \hline
13 & Borrar con el usuario de conexión & Permisos de conexión & 229 & permiso DELETE denegado \\ \hline
\end{longtable}
\end{center}


\captura{06-rechazos}{Captura 6. Las pruebas de rechazo en la consola.}

La base rechazó las \dnNumPruebasRechazadas{} operaciones, y cada una con la restricción que corresponde a su regla. El error 547 es el de una restricción \at{CHECK} o de una llave foránea; el 2627, el de una restricción \at{UNIQUE}; el 2601, el de un índice único filtrado, y el 229, el de un permiso denegado. Tres resultados merecen un comentario:

\begin{reglas}
  \item \textbf{Prueba 2.} \at{LUIS\_QUÓ} se considera el mismo nickname que \at{luis\_quo}, porque la columna compara sin distinguir mayúsculas ni acentos (\mbox{CU-02} \mbox{RN-01}).
  \item \textbf{Pruebas 6 y 7.} Las dos las rechazan las llaves foráneas compuestas de \textit{Equipamiento}: la primera, porque pedro no posee el objeto; la segunda, porque el objeto es un muro y la ranura es la del peón.
  \item \textbf{Prueba 13.} El usuario de conexión del servidor no puede borrar, ni siquiera un silencio, que es de las pocas filas que los casos de uso eliminan: para eso tiene los procedimientos almacenados.
\end{reglas}
